%% =====================================================================
%%  Capítulo 5 -- Estabilidad
%%  Migrado al formato nuevo (noviembre de 2026) desde Handout4_Ctrl.tex,
%%  con material de las Actividades 1 (Routh y Kharitonov) y de la clase 2
%%  de Bode (rango de Kc por Routh y por Bode).
%%  Cambios de fondo respecto a la versión anterior (ver MIGRACION.md):
%%    la sección de tiempo muerto (Padé) se conserva, con la fórmula de c_k
%%    corregida y la comparación con el límite exacto; "estable si todos los polos Y CEROS están
%%    en el semiplano izquierdo" -> solo los polos; el ejemplo
%%    s^3+2s^2+s+2 es un caso de FILA de ceros (no de epsilon); se
%%    resolvieron los ejemplos que no tenían solución; en el primer ejemplo
%%    de Kharitonov faltaba el intervalo de a4 y la familia NO es
%%    robustamente estable; el ejemplo de diseño robusto se resolvió y se
%%    muestra el conservadurismo; C(s)/R(s) -> Y(s)/R(s); Kharitanov ->
%%    Kharitonov; anécdota del sirviente (es de Kelvin, no de Maxwell).
%% =====================================================================

\capitulo[cap:es]{Estabilidad}
  {En la cocina, la educación y el amor la rapidez va en contra del placer y la calidad. Necesitamos cocinar, amar y leer a fuego lento.}
  {Julián de Zubiría Samper}

\begin{objetivos}[partes/05-estabilidad.html]
  \item Distinguir estabilidad asintótica, marginal (en el sentido de Lyapunov) y BIBO, y decidirla a partir de los polos o de los valores propios.
  \item Aplicar el criterio de Routh-Hurwitz, incluidos sus dos casos especiales, para contar las raíces con parte real positiva sin calcularlas.
  \item Usar la tabla de Routh para diseñar: rango de ganancia, ganancia y frecuencia críticas, regiones de estabilidad de un PI y estabilidad relativa.
  \item Analizar la estabilidad de un lazo con tiempo muerto mediante la aproximación de Padé, y reconocer cuánto se equivoca.
  \item Decidir la estabilidad de un equilibrio en el espacio de estado, también cuando el modelo viene de linealizar, y reconocer los modos que la función de transferencia esconde.
  \item Verificar la estabilidad robusta de una familia de polinomios de intervalo con el teorema de Kharitonov, y reconocer cuándo el resultado es conservador.
\end{objetivos}

%% ---------------------------------------------------------------------
\section{¿Qué es estabilidad?}
\label{sec:es-concepto}
%% ---------------------------------------------------------------------

Hasta aquí hemos analizado la respuesta transitoria: qué tan rápido responde un sistema y cuánto se pasa. Antes de preguntar \emph{cómo} llega, hay que preguntar \emph{si} llega. Esa es la pregunta de la estabilidad, y es la primera exigencia de cualquier diseño: un controlador que desestabiliza el lazo no sirve, por buenos que sean los demás números.

El diccionario de la Real Academia da varias acepciones de \emph{estable}; la que nos sirve es ``que mantiene o recupera el equilibrio''. La palabra clave es \clave{equilibrio}: la estabilidad no es una propiedad del sistema sino de un \alerta{punto de equilibrio} del sistema. Un péndulo tiene dos: colgado hacia abajo, donde una perturbación pequeña se amortigua, y parado hacia arriba, donde la menor perturbación lo tumba. Usted mismo, de pie, está en un equilibrio inestable que su sistema nervioso estabiliza todo el tiempo.

En un sistema lineal e invariante no forzado hay un solo equilibrio, el origen ($x=0$); con entrada constante, el equilibrio se desplaza, pero en variables desviación sigue siendo el origen. Como hay uno solo, se abusa del lenguaje y se dice ``el sistema es estable''. Lo haremos también, pero recuerde que siempre se habla del equilibrio.\nota{En sistemas no lineales el abuso no se puede cometer: un mismo sistema tiene equilibrios estables e inestables, como el péndulo (\cref{sec:es-espacio}).}

\begin{definicion}[label=def:es-estabilidad]{Tres nociones de estabilidad}
Para un sistema lineal e invariante, con su equilibrio en el origen:
\begin{itemize}
  \item El equilibrio es \clave{asintóticamente estable} si la respuesta natural (sin entrada, desde cualquier condición inicial) tiende a cero.
  \item Es \clave{marginalmente estable} (estable en el sentido de Lyapunov) si la respuesta natural se mantiene acotada, pero no necesariamente tiende a cero.
  \item Es \clave{inestable} si alguna condición inicial produce una respuesta natural que crece sin límite.
\end{itemize}
El sistema es \clave{BIBO estable} (\emph{bounded input, bounded output}) si toda entrada acotada produce una salida acotada.
\end{definicion}

\subsection*{Los polos lo deciden todo}

La respuesta natural es una suma de \clave{modos} $e^{p_it}$, uno por cada polo $p_i$ (capítulo~2). Un polo real negativo da un modo que decae; un par $\sigma\pm j\omega$ da $e^{\sigma t}\sin(\omega t+\phi)$, que decae si $\sigma<0$, oscila con amplitud constante si $\sigma=0$ y crece si $\sigma>0$. Los polos repetidos agregan factores $t,t^2,\dots$: sobre el eje imaginario, un polo doble da $t\sin(\omega t)$, que crece aunque no haya polos en el semiplano derecho (\cref{fig:es-modos}).

\begin{figure}[ht]
\centering
\includegraphics[width=\linewidth]{cap05/f_modos}
\caption{Ubicación de los polos y respuesta natural (al impulso). Solo la primera tiende a cero. La última muestra que los polos repetidos sobre el eje imaginario también desestabilizan.}
\label{fig:es-modos}
\end{figure}

\begin{ideaclave}
\textbf{Estabilidad y ubicación de los polos.}\enspace Sea $\Tlc(s)=N(s)/D(s)$ sin cancelaciones. El equilibrio es
\begin{itemize}
  \item \textbf{asintóticamente estable} si y solo si todos los polos tienen parte real negativa (semiplano izquierdo abierto). Solo en este caso el sistema es además BIBO estable;
  \item \textbf{marginalmente estable} si no hay polos en el semiplano derecho y los que están sobre el eje imaginario son simples;
  \item \textbf{inestable} si hay algún polo con parte real positiva o algún polo repetido sobre el eje imaginario.
\end{itemize}
Los \alerta{ceros no intervienen}: modifican la forma de la respuesta, no su estabilidad.
\end{ideaclave}

\begin{cuidado}[Marginal no es BIBO]
Un sistema marginalmente estable \emph{no} es BIBO estable. El integrador $1/s$ ante un escalón da una rampa, y $1/(s^2+1)$ ante $\sin t$ da $\tfrac12(\sin t-t\cos t)$, que crece como $t/2$: la entrada excita justo el modo que no se amortigua. En diseño, ``estable'' quiere decir siempre \emph{asintóticamente} estable.
\end{cuidado}

\codigorepo{cap05_estabilidad/C5_Ej1_Modos.m}{modos}{\cref{fig:es-modos}: la respuesta natural según la ubicación de los polos}

\subsection*{¿Cómo se decide?}

Para saber si un equilibrio es estable hay varios caminos, y el curso los recorre casi todos:
\begin{itemize}
  \item \textbf{Calcular los polos} o los valores propios (\cmd{roots}, \cmd{pole}, \cmd{eig}). Con computador es lo más directo, pero no dice \emph{para qué valores de un parámetro} el lazo es estable.
  \item \textbf{El criterio de Routh-Hurwitz} (\cref{sec:es-routh}): cuenta las raíces inestables a partir de los coeficientes, sin calcularlas. Como trabaja con los coeficientes, admite parámetros libres: es una herramienta de diseño.
  \item \textbf{El lugar de las raíces} de Evans~\citep{Eva50}, que dibuja cómo se mueven los polos del lazo cerrado cuando cambia una ganancia.
  \item \textbf{Los márgenes de fase y de ganancia} sobre el diagrama de Bode y el criterio de Nyquist~\citep{Nyq32,Bod40}, en el capítulo~10.
  \item \textbf{El método de Lyapunov}~\citep{Lya92}, el más general, que sirve también para sistemas no lineales y variantes en el tiempo; en este curso usaremos su método indirecto, la linealización (\cref{sec:es-espacio}).
\end{itemize}

\begin{pausa}
Clasifique: (a) $\dfrac{s-3}{(s+1)(s+4)}$; (b) $\dfrac{1}{s(s+2)}$; (c) $\dfrac{s+1}{(s^2+4)^2}$; (d) $\dfrac{1}{s^2-s+1}$.
\respuesta{(a) Asintóticamente estable: el cero en $+3$ no importa. (b) Marginalmente estable: polo simple en el origen (un escalón da una rampa: no es BIBO). (c) Inestable: polos dobles en $\pm 2j$. (d) Inestable: polos en $0.5\pm j0.866$.}
\end{pausa}

%% ---------------------------------------------------------------------
\section{Un poco de historia: de los anillos de Saturno a las turbinas de Davos}
\label{sec:es-historia}
%% ---------------------------------------------------------------------

En el capítulo~1 vimos el recorrido general; aquí nos detenemos en la historia de la pregunta que resuelve este capítulo.

\begin{figure}[ht]
\centering
\includegraphics[width=0.9\linewidth]{cap05/personajes}
\caption{De izquierda a derecha: James Clerk Maxwell (1831--1879), Edward John Routh (1831--1907), Aurel Stodola (1859--1942) y Adolf Hurwitz (1859--1919).}
\label{fig:es-personajes}
\end{figure}

\paragraph{Maxwell deja la pregunta abierta.} El primer análisis matemático de la estabilidad de un regulador fue el de George Biddell Airy, astrónomo real, que en 1840 estudió por qué el regulador de fricción de un telescopio oscilaba~\citep{Air40}. Pero la pregunta general la planteó James Clerk Maxwell en \emph{On governors} (1868)~\citep{Max68,Kan16}. Maxwell ya había usado la estabilidad como criterio de verdad: en su ensayo sobre los anillos de Saturno, premiado en 1857, descartó los modelos que no eran estables, porque unos anillos que existen desde hace tanto tiempo tienen que serlo~\citep{Bit15}. En \emph{On governors} linealizó las ecuaciones del regulador y de la máquina y mostró que el equilibrio es estable si todas las raíces de la ecuación característica tienen parte real negativa. Obtuvo las condiciones para segundo y tercer orden ---la de tercer orden es la que veremos en el ejemplo~\ref{ej:es-tercer}--- y confesó que no había podido hacerlo para grados mayores, con la esperanza de que el problema ``atrajera la atención de los matemáticos''.

\paragraph{Routh y el premio Adams.} La Universidad de Cambridge entregaba cada dos años el premio Adams, creado en honor de John Couch Adams, quien había predicho la existencia de Neptuno a partir de las irregularidades de la órbita de Urano. Para el periodo 1875--1877 los examinadores ---entre ellos Maxwell y George Stokes--- escogieron como tema \emph{El criterio de estabilidad dinámica}. Lo ganó en 1877 Edward John Routh, con el tratado que contiene el arreglo que hoy lleva su nombre~\citep{Rou77}. Routh y Maxwell habían nacido el mismo año, entraron juntos a Cambridge y presentaron juntos los exámenes de matemáticas de 1854, en los que Routh quedó primero y Maxwell segundo. Routh fue además el tutor más famoso de Cambridge: preparó a más de seiscientos estudiantes, entre ellos Rayleigh y J.~J.\ Thomson. Estaba casado con una hija de Airy.

\paragraph{Stodola, Hurwitz y las turbinas.} El trabajo de Routh tardó en salir de Inglaterra. En el continente el texto de referencia sobre reguladores era el de Ivan Vyshnegradskii (1876), que había resuelto el caso de tercer orden con un método gráfico~\citep{Vys76}. Sobre esa base trabajaba Aurel Stodola, ingeniero eslovaco que llegó en 1892 como profesor al Politécnico de Zúrich (hoy ETH). En 1893 le encargaron el regulador de las turbinas hidráulicas de una planta en Davos; su modelo era de orden mayor que tres, el criterio de Vyshnegradskii no alcanzaba, y Stodola le preguntó a su colega matemático Adolf Hurwitz cuándo un polinomio tiene todas sus raíces con parte real negativa~\citep{Bis89}. Hurwitz respondió con un criterio en términos de determinantes, publicado en 1895~\citep{Hur95}, sin conocer el de Routh. Stodola publicó su análisis de las turbinas en la revista suiza de ingeniería (1893--94) y presentó los resultados en variables relativas, sin dimensiones, lo que los hizo útiles para otros ingenieros.

\paragraph{Dos criterios, uno solo.} En 1911 el italiano Enrico Bompiani demostró que los criterios de Routh y de Hurwitz son equivalentes~\citep{Bom11}; ese mismo año Luciano Orlando publicó otra demostración y una fórmula que relaciona los determinantes de Hurwitz~\citep{Orl11}. Por eso hoy se habla del criterio de \clave{Routh-Hurwitz}: el arreglo de Routh es la forma cómoda de calcular, y los determinantes de Hurwitz son la forma cómoda de demostrar.

\begin{historia}[Lyapunov y Kharitonov]
En 1892 Aleksandr Lyapunov, entonces profesor en Járkov, presentó su tesis \emph{El problema general de la estabilidad del movimiento}~\citep{Lya92}. Definió con precisión qué es un equilibrio estable y dio dos métodos: el \emph{indirecto}, que decide la estabilidad local linealizando (\cref{sec:es-espacio}), y el \emph{directo}, que usa funciones parecidas a la energía y sirve para sistemas no lineales. Su obra se conoció en Occidente a mediados del siglo~XX. Casi un siglo después, en 1978, Vladimir Kharitonov publicó en una revista soviética de ecuaciones diferenciales un teorema sorprendente: para saber si \emph{toda} una familia de polinomios con coeficientes inciertos es estable, basta revisar cuatro polinomios~\citep{Kha78}. El resultado pasó inadvertido fuera de la Unión Soviética hasta 1983--84, cuando Ross Barmish y otros lo dieron a conocer en las revistas de control occidentales~\citep{Bar84} (\cref{sec:es-kharitonov}).
\end{historia}

%% ---------------------------------------------------------------------
\section{El criterio de Routh-Hurwitz}
\label{sec:es-routh}
%% ---------------------------------------------------------------------

Para el lazo de la \cref{fig:es-lazo}, la función de transferencia en lazo cerrado es
\[
  \frac{Y(s)}{R(s)}=\frac{\Gc(s)\Gp(s)}{1+\Gc(s)\Gp(s)\Hs(s)}
  =\frac{b_0s^m+b_1s^{m-1}+\cdots+b_m}{a_0s^n+a_1s^{n-1}+\cdots+a_{n-1}s+a_n},\qquad m\le n .
\]
El denominador igualado a cero es la \clave{ecuación característica},
\begin{equation}
  1+\Gma(s)=0\quad\Longleftrightarrow\quad a_0s^n+a_1s^{n-1}+\cdots+a_{n-1}s+a_n=0,
  \label{eq:es-caracteristica}
\end{equation}
con $\Gma=\Gc\Gp\Hs$ la malla abierta. Sus raíces son los polos del lazo cerrado: \alerta{la estabilidad del lazo se decide en el denominador de $1+\Gma$}, que es numerador más denominador de $\Gma$ cuando esta es una fracción sin cancelaciones.

\begin{figure}[ht]
\centering
\includegraphics[width=0.62\linewidth]{cap10/lazo_basico}
\caption{Lazo básico del curso.}
\label{fig:es-lazo}
\end{figure}

\subsection*{Una condición necesaria: los signos}

Si todas las raíces tienen parte real negativa, el polinomio es un producto de factores $(s+p)$ y $(s^2+2\zeta\wn s+\wn^2)$ con coeficientes positivos, y al multiplicarlos no aparece ningún signo menos ni ningún término faltante. Entonces:

\begin{ideaclave}
\textbf{Condición necesaria.}\enspace Si a la ecuación característica le \alerta{falta un término} o tiene \alerta{coeficientes de signos distintos}, hay raíces con parte real positiva o nula: el equilibrio no es asintóticamente estable. Para $n=1$ y $n=2$ la condición también es suficiente; para $n\ge3$ \emph{no}: $s^3+s^2+2s+8$ tiene todos sus coeficientes positivos y raíces en $0.5\pm j1.94$.
\end{ideaclave}

\subsection*{La tabla de Routh}

Con todos los coeficientes del mismo signo (tomamos $a_0>0$), se arma la tabla. Las dos primeras filas se llenan alternando los coeficientes; cada fila siguiente se calcula con las dos de arriba:
\[
\begin{array}{c|cccc}
  s^n     & a_0 & a_2 & a_4 & \cdots\\
  s^{n-1} & a_1 & a_3 & a_5 & \cdots\\
  s^{n-2} & b_1 & b_2 & b_3 & \cdots\\
  s^{n-3} & c_1 & c_2 & \cdots\\
  \vdots  & \vdots\\
  s^0     & \ast
\end{array}
\qquad
\begin{aligned}
  b_1&=\frac{a_1a_2-a_0a_3}{a_1}, & b_2&=\frac{a_1a_4-a_0a_5}{a_1}, & \dots\\[3pt]
  c_1&=\frac{b_1a_3-a_1b_2}{b_1}, & c_2&=\frac{b_1a_5-a_1b_3}{b_1}, & \dots
\end{aligned}
\]
Todas las entradas siguen el mismo patrón: ``el de abajo a la izquierda por el de arriba a la derecha, menos el de arriba a la izquierda por el de abajo a la derecha, dividido por el de abajo a la izquierda''. Donde falte un elemento se pone cero. La tabla tiene $n+1$ filas, de $s^n$ a $s^0$, y la última fila es siempre $a_n$.

\begin{definicion}[label=def:es-routh]{Criterio de Routh-Hurwitz}
El número de raíces de la ecuación característica con \alerta{parte real positiva} es igual al número de \clave{cambios de signo} en la primera columna de la tabla de Routh. El equilibrio es asintóticamente estable si y solo si todos los elementos de la primera columna tienen el mismo signo~\citep{Rou77,Hur95}.
\end{definicion}

\begin{cuidado}[Multiplicar una fila]
Una fila completa se puede multiplicar o dividir por un número \emph{positivo} sin cambiar los signos de la primera columna; sirve para evitar fracciones. Nunca por un número negativo, ni por una expresión cuyo signo no conoce (como $K$).
\end{cuidado}

\begin{ejemplo}[label=ej:es-tercer]{Tercer orden: la condición de Maxwell}
Encuentre la condición de estabilidad de $a_0s^3+a_1s^2+a_2s+a_3=0$, con $a_0>0$.
\solucion
\[
\begin{array}{c|cc}
  s^3 & a_0 & a_2\\
  s^2 & a_1 & a_3\\
  s^1 & \dfrac{a_1a_2-a_0a_3}{a_1} & 0\\
  s^0 & a_3
\end{array}
\]
Sin cambios de signo: $a_1>0$, $a_3>0$ y
\begin{equation}
  a_1a_2>a_0a_3 .
  \label{eq:es-tercer}
\end{equation}
Es la condición que encontró Maxwell en 1868. Vale la pena memorizarla: \alerta{``el producto de los del medio mayor que el producto de los extremos''}.
\end{ejemplo}

\begin{ejemplo}[label=ej:es-cuarto]{Cuarto orden}
¿Cuántas raíces inestables tiene $s^4+2s^3+3s^2+4s+5=0$?
\solucion Todos los coeficientes son positivos, así que hay que hacer la tabla:
\[
\begin{array}{c|ccc}
  s^4 & 1 & 3 & 5\\
  s^3 & 2 & 4 & 0\\
  s^2 & \frac{2\cdot3-1\cdot4}{2}=1 & \frac{2\cdot5-1\cdot0}{2}=5\\
  s^1 & \frac{1\cdot4-2\cdot5}{1}=-6 & 0\\
  s^0 & 5
\end{array}
\]
La primera columna es $1,\,2,\,1,\,-6,\,5$: dos cambios de signo ($1\to-6$ y $-6\to5$), luego \alerta{dos raíces con parte real positiva}. En efecto, \cmd{roots([1 2 3 4 5])} da $-1.288\pm j0.858$ y $0.288\pm j1.416$.
\end{ejemplo}

\codigorepo{cap05_estabilidad/C5_Ej2_Routh.m}{routh}{ejemplo~\ref{ej:es-cuarto} y el contraejemplo de los signos con la función \cmd{tablaRouth}}

\begin{receta}[Criterio de Routh-Hurwitz]
\begin{enumerate}
  \item Escribir la ecuación característica $1+\Gma(s)=0$ como polinomio, ordenado en potencias decrecientes.
  \item Revisar los signos: si falta un término o hay signos distintos, ya se sabe que no es asintóticamente estable (la tabla dirá cuántas raíces son inestables).
  \item Armar la tabla. Si aparece un cero en la primera columna o una fila de ceros, ir a la \cref{sec:es-casos}.
  \item Contar los cambios de signo de la primera columna.
  \item Verificar con \cmd{roots}, \cmd{pole} o \cmd{tablaRouth}.
\end{enumerate}
\end{receta}

\begin{pausa}
Use \eqref{eq:es-tercer} para decidir, sin tabla, si $s^3+2s^2+3s+4$ y $2s^3+s^2+3s+4$ son estables.
\respuesta{El primero sí: $2\cdot3=6>1\cdot4$. El segundo no: $1\cdot3=3<2\cdot4=8$; la tabla tiene dos cambios de signo y dos raíces inestables.}
\end{pausa}

%% ---------------------------------------------------------------------
\section{Los dos casos especiales}
\label{sec:es-casos}
%% ---------------------------------------------------------------------

La regla para construir la tabla divide por el primer elemento de la fila anterior. Si ese elemento es cero, hay que hacer algo. Hay dos situaciones distintas, y conviene no confundirlas.

\subsection*{Caso 1: un cero en la primera columna, con el resto de la fila distinto de cero}

Se reemplaza el cero por un número pequeño y positivo $\epsilon$, se sigue construyendo la tabla y se cuentan los cambios de signo haciendo $\epsilon\to0^+$.

\begin{ejemplo}[label=ej:es-caso1]{Un cero que se vuelve $\epsilon$}
Analice $s^5+2s^4+4s^3+8s^2+10s+6=0$.
\solucion
\[
\begin{array}{c|ccc}
  s^5 & 1 & 4 & 10\\
  s^4 & 2 & 8 & 6\\
  s^3 & 0\;\to\;\epsilon & 7 & 0\\
  s^2 & 8-\dfrac{14}{\epsilon} & 6\\[6pt]
  s^1 & \approx 7 & \\
  s^0 & 6
\end{array}
\]
La fila $s^3$ vale $\bigl[\frac{2\cdot4-1\cdot8}{2},\,\frac{2\cdot10-1\cdot6}{2}\bigr]=[0,\,7]$: solo el primer elemento es cero. Con $\epsilon\to0^+$, el elemento de $s^2$ es muy negativo y el de $s^1$, $\dfrac{(8-14/\epsilon)\,7-6\epsilon}{8-14/\epsilon}$, tiende a $7$. La primera columna queda $+,+,+,-,+,+$: \alerta{dos cambios de signo}, dos raíces inestables ($0.484\pm j1.738$).
\end{ejemplo}

\subsection*{Caso 2: una fila completa de ceros}

Una fila entera se anula cuando el polinomio tiene un factor \clave{par} (solo potencias pares de $s$), cuyas raíces son simétricas respecto al origen (\cref{fig:es-simetricas}): reales y opuestas (A), imaginarias puras (B) o en cuadrante, $\pm\sigma\pm j\omega$ (C). Ese factor es el \clave{polinomio auxiliar} $P(s)$, que se lee en la fila \emph{anterior} a la de ceros. Se reemplaza la fila de ceros por los coeficientes de $\dif P/\dif s$ y se sigue.

\begin{figure}[ht]
\centering
\includegraphics[width=0.36\linewidth]{cap05/raices_simetricas}
\caption{Raíces simétricas respecto al origen que producen una fila de ceros: reales (A), imaginarias (B) o en cuadrante (C).}
\label{fig:es-simetricas}
\end{figure}

\begin{ideaclave}
\textbf{Cómo leer el caso 2.}\enspace \begin{itemize}
  \item Las raíces del polinomio auxiliar son raíces del polinomio original, y se pueden calcular: $P(s)$ tiene grado par y solo potencias pares.
  \item Los cambios de signo \emph{desde la fila del auxiliar hacia abajo} cuentan las raíces del auxiliar con parte real positiva; como son simétricas, hay el mismo número con parte real negativa, y las que faltan están sobre el eje imaginario.
  \item Si no hay cambios de signo y el auxiliar tiene raíces imaginarias simples, el equilibrio es \alerta{marginalmente estable}.
\end{itemize}
\end{ideaclave}

\begin{ejemplo}[label=ej:es-caso2]{Raíces sobre el eje imaginario}
Analice $s^5+7s^4+6s^3+42s^2+8s+56=0$.
\solucion
\[
\begin{array}{c|ccc}
  s^5 & 1 & 6 & 8\\
  s^4 & 7 & 42 & 56\\
  s^3 & 0\;\to\;28 & 0\;\to\;84 & \\
  s^2 & 21 & 56\\
  s^1 & 28/3\\
  s^0 & 56
\end{array}
\]
La fila $s^3$ es toda cero: $\frac{7\cdot6-1\cdot42}{7}=0$ y $\frac{7\cdot8-1\cdot56}{7}=0$. El auxiliar es la fila $s^4$: $P(s)=7s^4+42s^2+56=7(s^2+2)(s^2+4)$, y su derivada $28s^3+84s$ da la nueva fila. No hay cambios de signo, así que no hay raíces en el semiplano derecho; las cuatro raíces del auxiliar están sobre el eje: $\pm j\sqrt2$ y $\pm 2j$. La quinta es $-7$. El equilibrio es \alerta{marginalmente estable}: la respuesta natural oscila con dos frecuencias sin amortiguarse.
\end{ejemplo}

\begin{ejemplo}[label=ej:es-caso2b]{Un caso 2 que parece caso 1}
Analice $s^3+2s^2+s+2=0$.
\solucion La fila $s^1$ vale $\frac{2\cdot1-1\cdot2}{2}=0$, y como es su único elemento, es una \emph{fila completa} de ceros: caso 2, no caso 1. El auxiliar es $2s^2+2$, con derivada $4s$:
\[
\begin{array}{c|cc}
  s^3 & 1 & 1\\
  s^2 & 2 & 2\\
  s^1 & 0\;\to\;4\\
  s^0 & 2
\end{array}
\]
Sin cambios de signo y auxiliar $2(s^2+1)$: raíces $\pm j$ y $-2$. Marginalmente estable. (Reemplazar el cero por $\epsilon$ da los mismos signos, pero esconde la información importante: dónde están las raíces del eje y a qué frecuencia oscilará el sistema.)
\end{ejemplo}

\codigorepo{cap05_estabilidad/C5_Ej3_CasosEspeciales.m}{caso2}{ejemplo~\ref{ej:es-caso2}: el polinomio auxiliar y sus raíces}

\begin{pausa}
Analice $s^4+4=0$. Antes de hacer la tabla, ¿qué puede decir solo con mirar los coeficientes?
\respuesta{Faltan los términos en $s^3$, $s^2$ y $s$: no es asintóticamente estable. Toda la fila $s^3$ es cero; el auxiliar es $s^4+4$, con derivada $4s^3$. La fila $s^2$ queda $[0,\,4]$ (caso 1: $\epsilon$) y la fila $s^1$, $-16/\epsilon$. Dos cambios de signo: dos raíces en el semiplano derecho. Son $\pm1\pm j$, en cuadrante (tipo C).}
\end{pausa}

%% ---------------------------------------------------------------------
\section{Routh-Hurwitz para diseñar}
\label{sec:es-diseno}
%% ---------------------------------------------------------------------

La gran ventaja de Routh-Hurwitz sobre \cmd{roots} es que trabaja con los coeficientes, y los coeficientes pueden contener parámetros del controlador. La tabla responde entonces la pregunta de diseño más básica: \emph{¿para qué valores de la ganancia el lazo es estable?}

\subsection{Rango de ganancia, ganancia crítica y frecuencia de oscilación}

\begin{ejemplo}[label=ej:es-rangokc]{Rango de $\Kc$ de un lazo de proceso}
En el lazo de la \cref{fig:es-tres}, un controlador proporcional $\Kc$ actúa sobre una válvula, un proceso y un sensor de primer orden. Halle el rango de $\Kc$ para el cual el lazo es estable, la ganancia crítica y la frecuencia a la que oscila.
\solucion \textbf{1. Ecuación característica.}
\[
  1+\Kc\,\frac{0.016}{3s+1}\cdot\frac{50}{30s+1}\cdot\frac{1}{10s+1}=0
  \;\Longrightarrow\;
  900s^3+420s^2+43s+(1+0.8\Kc)=0 .
\]
\textbf{2. Tabla.}
\[
\begin{array}{c|cc}
  s^3 & 900 & 43\\
  s^2 & 420 & 1+0.8\Kc\\
  s^1 & \dfrac{17\,160-720\Kc}{420} & \\[6pt]
  s^0 & 1+0.8\Kc
\end{array}
\]
\textbf{3. Condiciones.} Fila $s^0$: $1+0.8\Kc>0$, o sea $\Kc>-1.25$. Fila $s^1$: $17\,160-720\Kc>0$, o sea $\Kc<23.83$. Entonces
\[
  \alerta{-1.25<\Kc<23.83}.
\]
\textbf{4. Ganancia y frecuencia críticas.} En $\Kc=\Kcr=23.83$ la fila $s^1$ se anula: es una fila de ceros, y el auxiliar es la fila $s^2$,
\[
  420s^2+(1+0.8\cdot23.83)=420s^2+20.07=0\;\Longrightarrow\;s=\pm j\,0.2186 .
\]
El lazo oscila sin amortiguarse con $\wcr=0.2186$\,rad/s, es decir, con periodo $\Pcr=2\pi/\wcr=28.7$\,s (\cref{fig:es-rangokc}).
\end{ejemplo}

\begin{figure}[ht]
\centering
\includegraphics[width=0.92\linewidth]{cap05/lazo_tres}
\caption{Lazo de proceso del ejemplo~\ref{ej:es-rangokc}.}
\label{fig:es-tres}
\end{figure}

\begin{figure}[ht]
\centering
\includegraphics[width=0.85\linewidth]{cap05/f_rango_kc}
\caption{Ejemplo~\ref{ej:es-rangokc}: por debajo de $\Kcr$ la oscilación se amortigua, en $\Kcr$ se sostiene con periodo $\Pcr=28.7$\,s y por encima crece.}
\label{fig:es-rangokc}
\end{figure}

\begin{ideaclave}
\textbf{Ganancia crítica: tres caras del mismo número.}\enspace $\Kcr$ y $\wcr$ salen de la fila que se anula y del polinomio auxiliar. Son los mismos números que (i) el método de oscilación sostenida de Ziegler y Nichols busca experimentalmente en la planta (capítulo~8)~\citep{ZiNi42}, y (ii) el margen de ganancia y la frecuencia de cruce de fase que se leen en el diagrama de Bode (capítulo~10): para este lazo, \cmd{margin} da $\GM=23.83$ y $\wgm=0.2186$\,rad/s.
\end{ideaclave}

\codigorepo{cap05_estabilidad/C5_Ej4_RangoK.m}{rangokc}{ejemplo~\ref{ej:es-rangokc}: rango de $\Kc$, $\Kcr$, $\wcr$ y la comparación con \cmd{margin}}

\begin{cuidado}[Las ganancias negativas]
El rango incluye $-1.25<\Kc<0$. Con $\Kc$ negativo el controlador empuja en la dirección equivocada; el lazo resulta estable mientras el término independiente siga positivo, pero el error en estado estacionario es peor que en lazo abierto. Que un valor sea estable no lo hace útil: en la práctica se usa $0<\Kc<\Kcr$, y bastante por debajo de $\Kcr$.
\end{cuidado}

\subsection{Estabilizar una planta inestable}

Cuando la planta es inestable, la ganancia ya no tiene una cota superior sola: también aparece una cota \emph{inferior}, porque con poca ganancia el controlador no alcanza a contener el polo inestable.

\begin{ejemplo}[label=ej:es-inestable]{Una ventana estrecha}
La planta $\Gp(s)=\dfrac{1}{(s-1)(s^2+2s+2)}$ se controla con el compensador $\Gc(s)=\dfrac{K(s+2)}{s+3}$, con retroalimentación unitaria (\cref{fig:es-inestable-lazo}). Halle el rango de $K$ que estabiliza el lazo.
\solucion La ecuación característica es $(s+3)(s-1)(s^2+2s+2)+K(s+2)=0$:
\[
  s^4+4s^3+3s^2+(K-2)\,s+(2K-6)=0 .
\]
\[
\begin{array}{c|ccc}
  s^4 & 1 & 3 & 2K-6\\
  s^3 & 4 & K-2 & \\
  s^2 & \dfrac{14-K}{4} & 2K-6\\[6pt]
  s^1 & \dfrac{68-16K-K^2}{14-K} & \\[6pt]
  s^0 & 2K-6
\end{array}
\]
Las condiciones son $K<14$ (fila $s^2$), $2K-6>0$ (fila $s^0$) y $68-16K-K^2>0$ (fila $s^1$, con $14-K>0$). Las raíces de $K^2+16K-68$ son $-8\pm2\sqrt{33}$, es decir $3.489$ y $-19.49$. Juntando:
\[
  \alerta{3<K<3.489}.
\]
Es un rango muy estrecho (\cref{fig:es-inestable}): un error de pocos puntos porcentuales en la ganancia, o en el modelo, desestabiliza el lazo. Este compensador estabiliza, pero no es robusto; dentro del rango, la mejor ganancia es $K\approx3.07$, con los polos más lentos en $-0.28$.
\end{ejemplo}

\begin{figure}[ht]
\centering
\includegraphics[width=0.78\linewidth]{cap05/lazo_inestable}
\caption{Ejemplo~\ref{ej:es-inestable}: compensador para una planta con un polo en $+1$.}
\label{fig:es-inestable-lazo}
\end{figure}

\begin{figure}[ht]
\centering
\includegraphics[width=0.85\linewidth]{cap05/f_inestable}
\caption{Ejemplo~\ref{ej:es-inestable}: parte real del polo más a la derecha en función de $K$. Por debajo de $3$ queda un polo real inestable; por encima de $3.489$ cruzan el eje dos polos complejos.}
\label{fig:es-inestable}
\end{figure}

\subsection{Dos parámetros: la región de estabilidad de un PI}

Con dos parámetros, las condiciones de la tabla dibujan una \emph{región} en el plano de los parámetros.

\begin{ejemplo}[label=ej:es-pi]{Región de estabilidad de un PI}
En el lazo de la \cref{fig:es-pi-lazo}, $\Gc(s)=\Kc+K_i/s$ y $\Gp(s)=\dfrac{2}{s^3+4s^2+5s+2}=\dfrac{2}{(s+1)^2(s+2)}$. Halle la región del plano $(\Kc,K_i)$ en la que el lazo es estable.
\solucion La ecuación característica es $s(s^3+4s^2+5s+2)+2(\Kc s+K_i)=0$:
\[
  s^4+4s^3+5s^2+(2+2\Kc)\,s+2K_i=0 .
\]
\[
\begin{array}{c|ccc}
  s^4 & 1 & 5 & 2K_i\\
  s^3 & 4 & 2+2\Kc & \\
  s^2 & \dfrac{9-\Kc}{2} & 2K_i\\[6pt]
  s^1 & c_1 & \\
  s^0 & 2K_i
\end{array}
\qquad
c_1=(2+2\Kc)-\frac{16K_i}{9-\Kc}.
\]
Las condiciones son $K_i>0$, $\Kc<9$ y $c_1>0$:
\[
  \alerta{0<K_i<\frac{(1+\Kc)(9-\Kc)}{8},\qquad -1<\Kc<9}.
\]
La frontera es una parábola con máximo $K_i=25/8$ en $\Kc=4$ (\cref{fig:es-region}). Sobre ella, la fila $s^1$ se anula y el auxiliar $\frac{9-\Kc}{2}s^2+2K_i$ da la frecuencia de la oscilación: con $\Kc=5$ y $K_i=3$, $\omega=\sqrt{2\cdot3/2}=1.732$\,rad/s (periodo $3.63$\,s).
\end{ejemplo}

\begin{figure}[ht]
\centering
\includegraphics[width=0.85\linewidth]{cap05/lazo_pi_estab}
\caption{Lazo del ejemplo~\ref{ej:es-pi}, con el PI en forma paralela.}
\label{fig:es-pi-lazo}
\end{figure}

\begin{figure}[ht]
\centering
\includegraphics[width=\linewidth]{cap05/f_region_pi}
\caption{Ejemplo~\ref{ej:es-pi}. Izquierda: región de estabilidad. Derecha: con $\Kc=5$, $K_i=1$ (adentro) es estable, aunque con $\OS=51\%$; $K_i=3$ (frontera) oscila sin amortiguarse; $K_i=4$ (afuera) es inestable.}
\label{fig:es-region}
\end{figure}

\codigorepo{cap05_estabilidad/C5_Ej5_RegionPI.m}{region}{ejemplo~\ref{ej:es-pi}: la región de estabilidad en el plano $(\Kc,K_i)$}

\begin{modelo}[LazoPI.slx]
El script \cmd{crearModelo\_LazoPI.m} arma el lazo de la \cref{fig:es-pi-lazo} con $\Kc$ y $K_i$ como variables del espacio de trabajo. Recorra la región: un punto adentro, uno sobre la frontera y uno afuera, y mida el periodo de la oscilación sostenida.
\end{modelo}

\begin{cuidado}[Estable no quiere decir aceptable]
El punto $(\Kc,K_i)=(5,1)$ está dentro de la región y, sin embargo, la respuesta tiene un sobrepico del $51\%$ y tarda $22$\,s en establecerse. La tabla de Routh dice \emph{si} es estable; para saber \emph{qué tan} estable hace falta algo más.
\end{cuidado}

\subsection{Estabilidad relativa}

La tabla se puede usar también para exigir un margen: que todos los polos tengan parte real menor que $-\sigma$, no solo menor que cero. Basta con desplazar el eje: se reemplaza $s=z-\sigma$ en la ecuación característica y se aplica el criterio al polinomio en $z$. Las raíces con $\Re z>0$ son las que tienen $\Re s>-\sigma$.

\begin{ejemplo}[label=ej:es-relativa]{Polos a la izquierda de $-0.05$}
En el ejemplo~\ref{ej:es-rangokc}, ¿para qué valores de $\Kc$ todos los polos tienen parte real menor que $-0.05$ (constante de tiempo de la envolvente menor que $20$\,s)?
\solucion Con $s=z-0.05$, $900s^3+420s^2+43s+(1+0.8\Kc)$ se convierte en
\[
  900z^3+285z^2+7.75z+(0.8\Kc-0.2125)=0 .
\]
\[
\begin{array}{c|cc}
  z^3 & 900 & 7.75\\
  z^2 & 285 & 0.8\Kc-0.2125\\
  z^1 & \frac{160-48\Kc}{19} & \\
  z^0 & 0.8\Kc-0.2125
\end{array}
\qquad\Longrightarrow\qquad
\alerta{0.266<\Kc<3.33}.
\]
El margen exigido reduce mucho el rango: de $(-1.25,\,23.83)$ a $(0.266,\,3.33)$. En los dos extremos hay un polo exactamente en $-0.05$ (\cref{fig:es-relativa}).
\end{ejemplo}

\begin{figure}[ht]
\centering
\includegraphics[width=0.5\linewidth]{cap05/f_relativa}
\caption{Ejemplo~\ref{ej:es-relativa}: polos del lazo para $\Kc=0.266$ y $\Kc=3.33$ (extremos del rango con margen) y para $\Kc=\Kcr$, con polos sobre el eje imaginario.}
\label{fig:es-relativa}
\end{figure}

\begin{pausa}
Para $\Gma(s)=\dfrac{K}{s(s+1)(s+2)}$ con retroalimentación unitaria, halle el rango de $K$, $\Kcr$ y $\wcr$.
\respuesta{$s^3+3s^2+2s+K=0$. Por \eqref{eq:es-tercer}: $3\cdot2>K$, y $K>0$: $0<K<6$. En $\Kcr=6$ el auxiliar es $3s^2+6=0$: $\wcr=\sqrt2=1.414$\,rad/s.}
\end{pausa}

%% ---------------------------------------------------------------------
\section{Tiempo muerto: la aproximación de Padé}
\label{sec:es-retardo}
%% ---------------------------------------------------------------------

En muchos procesos la salida no responde de inmediato: el agua tarda en recorrer la tubería entre la válvula y el sensor, la banda transportadora tarda en llevar el material a la báscula, la señal tarda en viajar por la red. La salida es la entrada \emph{desplazada} un tiempo $T$, el \clave{tiempo muerto} o \clave{retardo de transporte}. Por la propiedad de desplazamiento en el tiempo de la transformada de Laplace (capítulo~2),
\[
  y(t)=u(t-T)\,\escalon(t-T)\quad\Longleftrightarrow\quad Y(s)=e^{-Ts}\,U(s).
\]
Este es el único lugar del curso en que se analiza la estabilidad con retardo, y lo volveremos a encontrar en los modelos de primer orden con tiempo muerto que se usan para sintonizar controladores PID (capítulo~8).

\begin{cuidado}[Routh no se puede aplicar directamente]
Con un retardo, la ecuación característica deja de ser un polinomio: para el lazo de la \cref{fig:es-retardo} es
\[
  s+1+Ke^{-Ts}=0 ,
\]
una ecuación \emph{trascendente}, con infinitas raíces. La tabla de Routh necesita coeficientes de un polinomio, así que primero hay que \emph{aproximar} $e^{-Ts}$ por un cociente de polinomios.
\end{cuidado}

\begin{figure}[ht]
\centering
\includegraphics[width=0.8\linewidth]{cap05/lazo_retardo}
\caption{Lazo con retardo de transporte del ejemplo~\ref{ej:es-retardo}.}
\label{fig:es-retardo}
\end{figure}

\subsection*{La aproximación de Padé}

La \clave{aproximación de Padé} de orden $n$ es el cociente de dos polinomios de grado $n$ cuya serie de Taylor coincide con la de $e^{-Ts}$ en tantos términos como es posible~\citep{Pad92}:
\begin{equation}
  e^{-Ts}\approx P_n(s)=\frac{\sum_{k=0}^{n}(-1)^k c_k\,(Ts)^k}{\sum_{k=0}^{n}c_k\,(Ts)^k},
  \qquad
  c_k=\frac{(2n-k)!\;n!}{(2n)!\;k!\;(n-k)!},\quad k=0,\dots,n.
  \label{eq:es-pade}
\end{equation}
Con $n=1$, $c_0=1$ y $c_1=\tfrac12$; con $n=2$, $c_0=1$, $c_1=\tfrac12$ y $c_2=\tfrac1{12}$:
\[
  e^{-Ts}\approx\frac{1-\frac{T}{2}s}{1+\frac{T}{2}s},
  \qquad
  e^{-Ts}\approx\frac{1-\frac{T}{2}s+\frac{(Ts)^2}{12}}{1+\frac{T}{2}s+\frac{(Ts)^2}{12}} .
\]
Observe la estructura: el numerador tiene los mismos coeficientes que el denominador con los signos alternados, así que tiene \emph{ceros en el semiplano derecho} (simétricos a los polos). Por eso la aproximación tiene magnitud $1$ en todas las frecuencias, como el retardo, y por eso su respuesta al escalón arranca en sentido contrario (\cref{fig:es-pade}, izquierda): es la forma que tiene una función racional de ``esperar'' $T$ segundos.

\begin{historia}[Henri Padé]
El matemático francés Henri Padé (1863--1953) estudió en su tesis doctoral de 1892 cómo aproximar una función por cocientes de polinomios, y organizó las aproximaciones en la tabla que lleva su nombre~\citep{Pad92}. La de $e^{-x}$ fue uno de sus ejemplos. Setenta años después, cuando los ingenieros de control quisieron llevar los retardos a las herramientas de los sistemas racionales (Routh, el lugar de las raíces, el espacio de estado), la tomaron de ahí; en \matlab\ es la función \cmd{pade(T,n)}.
\end{historia}

\begin{ejemplo}[label=ej:es-retardo]{Rango de $K$ con tiempo muerto}
Halle, en función de $T$, el rango de $K$ para el cual el lazo de la \cref{fig:es-retardo} tiene un equilibrio asintóticamente estable.
\solucion \textbf{1. Aproximar.} Con Padé de primer orden, $1+K\dfrac{1}{s+1}\cdot\dfrac{1-\frac{T}{2}s}{1+\frac{T}{2}s}=0$, es decir,
\[
  (s+1)\Bigl(1+\tfrac{T}{2}s\Bigr)+K\Bigl(1-\tfrac{T}{2}s\Bigr)=0
  \;\Longrightarrow\;
  \tfrac{T}{2}\,s^2+\Bigl(1+\tfrac{T}{2}-K\tfrac{T}{2}\Bigr)s+(1+K)=0 .
\]
\textbf{2. Tabla.}
\[
\begin{array}{c|cc}
  s^2 & \frac{T}{2} & 1+K\\[3pt]
  s^1 & 1+\frac{T}{2}-K\frac{T}{2} & \\[3pt]
  s^0 & 1+K
\end{array}
\]
\textbf{3. Condiciones.} $1+K>0$ y $1+\frac{T}{2}(1-K)>0$:
\[
  \alerta{-1<K<1+\frac{2}{T}} .
\]
El rango se \emph{estrecha} a medida que crece el retardo: con $T=1$\,s, $K<3$; con $T=5$\,s, $K<1.4$. Y si $T\to0$, la cota superior desaparece, como debe ser: sin retardo el lazo es de primer orden, estable para todo $K>-1$.
\end{ejemplo}

\begin{ideaclave}
\textbf{El retardo consume estabilidad.}\enspace El retardo no cambia la magnitud de la malla abierta, pero le resta una fase $\omega T$ que crece sin límite con la frecuencia. Por eso un lazo que sin retardo era estable para cualquier ganancia, con retardo tiene una ganancia crítica, y esta baja cuando $T$ crece. Es la misma idea del margen de fase del capítulo~10: el retardo se ``come'' el margen.
\end{ideaclave}

\subsection*{¿Qué tan buena es la aproximación?}

Para este lazo el límite exacto se puede calcular. En el borde de la estabilidad hay raíces en $s=\pm j\omega$; reemplazando en $1+Ke^{-j\omega T}/(j\omega+1)=0$ y separando magnitud y fase,
\[
  K=\sqrt{1+\omega^2},\qquad \omega T+\arctan\omega=\pi .
\]
Con $T=1$: $\omega=2.029$\,rad/s y $K_{\max}=2.262$. Padé de primer orden predice $3$: un error del $33\%$, y \alerta{hacia el lado peligroso}. Con Padé de segundo orden la tabla da $2.292$, apenas $1.3\%$ por encima (\cref{fig:es-pade}, centro).

\begin{figure}[ht]
\centering
\includegraphics[width=\linewidth]{cap05/f_pade}
\caption{Izquierda: escalón retrasado $1$\,s y sus aproximaciones de Padé (arrancan en sentido contrario). Centro: ganancia máxima en función del retardo; Padé~1 sobrestima entre $15$ y $33\%$, Padé~2 alrededor del $1\%$. Derecha: con $K=2.5$ y $T=1$, el modelo de Padé~1 es estable, pero el lazo con el retardo exacto oscila y crece.}
\label{fig:es-pade}
\end{figure}

\begin{cuidado}[Padé de primer orden es optimista]
La aproximación de primer orden sobrestima la ganancia máxima: un diseño que la tabla declara estable puede no serlo (\cref{fig:es-pade}, derecha). Úsela para entender la tendencia ($K_{\max}$ baja con $T$) y para un primer cálculo; para decidir, use Padé de segundo orden o más, o verifique con el retardo exacto: \cmd{margin} y \cmd{allmargin} aceptan modelos con retardo (\texttt{tf(1,[1 1],'InputDelay',T)}), y en \simulink\ está el bloque \emph{Transport Delay}.
\end{cuidado}

\codigorepo{cap05_estabilidad/C5_Ej8_Pade.m}{pade}{ejemplo~\ref{ej:es-retardo}: Padé de orden 1 y 2 contra el retardo exacto}

\begin{modelo}[Retardo.slx]
El script \cmd{crearModelo\_Retardo.m} arma el lazo de la \cref{fig:es-retardo} dos veces: una con el bloque \emph{Transport Delay} y otra con la aproximación de Padé de primer orden. Con $T=1$, suba $K$ desde $2$ hasta $3$ y encuentre la ganancia a la que cada uno empieza a oscilar.
\end{modelo}

\begin{pausa}
Para $\Gma(s)=\dfrac{Ke^{-Ts}}{s+2}$, use Padé de primer orden para hallar el rango de $K$ en función de $T$. ¿Cuánto vale la cota con $T=1$?
\respuesta{$(s+2)(1+\tfrac{T}{2}s)+K(1-\tfrac{T}{2}s)=\tfrac{T}{2}s^2+(1+T-K\tfrac{T}{2})s+(2+K)$: $-2<K<2+\tfrac{2}{T}$. Con $T=1$, $K<4$ (el valor exacto es $3.04$).}
\end{pausa}

%% ---------------------------------------------------------------------
\section{Estabilidad en el espacio de estado}
\label{sec:es-espacio}
%% ---------------------------------------------------------------------

Con el modelo de estado $\dot x=\Am x+\Bm u$, $y=\Cm x+\Dm u$ (capítulo~4), la respuesta natural es $x(t)=e^{\Am t}x(0)$, y sus modos son $e^{\lambda_it}$, con $\lambda_i$ los \clave{valores propios} de $\Am$: las raíces de
\[
  \det(s\Id-\Am)=0 .
\]
Como $\dfrac{Y(s)}{U(s)}=\Cm(s\Id-\Am)^{-1}\Bm+\Dm=\dfrac{\Cm\operatorname{adj}(s\Id-\Am)\Bm}{\det(s\Id-\Am)}+\Dm$, los polos de la función de transferencia están entre los valores propios de $\Am$.

\begin{definicion}[label=def:es-hurwitz]{Matriz Hurwitz}
Una matriz $\Am$ es \clave{Hurwitz} si todos sus valores propios tienen parte real negativa. El equilibrio $x=0$ de $\dot x=\Am x$ es asintóticamente estable si y solo si $\Am$ es Hurwitz. Para decidirlo sin calcular los valores propios, se aplica Routh-Hurwitz al polinomio característico $\det(s\Id-\Am)$.
\end{definicion}

\begin{ejemplo}[label=ej:es-matriz]{Routh sobre el polinomio característico}
¿Es Hurwitz $\Am=\begin{bmatrix}0&1&0\\0&0&1\\-6&-11&-6\end{bmatrix}$?
\solucion Está en variables de fase, así que el polinomio característico se lee en la última fila: $s^3+6s^2+11s+6$. Por \eqref{eq:es-tercer}, $6\cdot11=66>1\cdot6$: es Hurwitz. (Sus valores propios son $-1$, $-2$ y $-3$.)
\end{ejemplo}

\subsection*{Lo que la función de transferencia esconde}

``Los polos están entre los valores propios'' quiere decir que pueden ser \emph{menos}: si hay cancelaciones, la función de transferencia no muestra todos los modos (\cref{sec:re-kalman}). Y un modo que no se ve puede ser inestable.

\begin{ejemplo}[label=ej:es-escondido]{Un modo inestable escondido}
Sea $\Am=\begin{bmatrix}0&1&0\\0&0&1\\2&1&-2\end{bmatrix}$, $\Bm=\begin{bmatrix}0\\0\\1\end{bmatrix}$, $\Cm=[-1\;\;0\;\;1]$.
\solucion El polinomio característico es $s^3+2s^2-s-2=(s-1)(s+1)(s+2)$: hay un signo menos, y el valor propio $+1$ hace que el equilibrio sea \alerta{inestable}. Sin embargo, la función de transferencia es
\[
  \frac{Y(s)}{U(s)}=\frac{s^2-1}{(s-1)(s+1)(s+2)}=\frac{1}{s+2},
\]
que parece perfectamente estable. El modo $e^{t}$ está ahí: cualquier condición inicial que lo excite hace crecer los estados, aunque no se vea en la salida.
\end{ejemplo}

\begin{cuidado}[Routh sobre $\det(s\Id-\Am)$, no sobre la función de transferencia]
Para decidir la estabilidad del \emph{equilibrio} se analiza el polinomio característico completo de $\Am$. Analizar el denominador de una función de transferencia ya simplificada solo dice si el sistema es BIBO estable \emph{entre esa entrada y esa salida}. Por la misma razón, nunca se cancela un polo inestable con un cero del controlador.
\end{cuidado}

\codigorepo{cap05_estabilidad/C5_Ej6_EspacioEstado.m}{escondido}{ejemplo~\ref{ej:es-escondido}: el modo inestable que la función de transferencia no muestra}

\subsection*{Equilibrios de sistemas no lineales: el método indirecto de Lyapunov}

Un sistema no lineal $\dot x=f(x,u)$ puede tener varios equilibrios. Para cada uno, la linealización de la \cref{sec:so-linealizacion} da $\dot{\dev x}=\Am\,\dev x+\Bm\,\dev u$, con $\Am=\partial f/\partial x$ evaluada en el equilibrio. El \clave{método indirecto de Lyapunov} dice:
\begin{itemize}
  \item si $\Am$ es Hurwitz, el equilibrio del sistema no lineal es \emph{localmente} asintóticamente estable;
  \item si $\Am$ tiene algún valor propio con parte real positiva, el equilibrio es inestable;
  \item si los valores propios con parte real máxima están sobre el eje imaginario, la linealización \alerta{no decide}: hay que mirar los términos no lineales.
\end{itemize}

\begin{ejemplo}[label=ej:es-pendulo]{Los dos equilibrios del péndulo}
El péndulo con fricción $\ddot\theta=-\frac{g}{l}\sin\theta-b\,\dot\theta$, con $g/l=9.81$\,s$^{-2}$ y $b=0.5$\,s$^{-1}$, tiene equilibrios en $\theta^*=0$ (abajo) y $\theta^*=\pi$ (arriba). Decida su estabilidad.
\solucion Con $x_1=\theta$, $x_2=\dot\theta$, el jacobiano es $\Am=\begin{bmatrix}0&1\\-\frac{g}{l}\cos\theta^*&-b\end{bmatrix}$.
\begin{itemize}
  \item Abajo, $\cos0=1$: $s^2+0.5s+9.81$, coeficientes positivos y segundo orden: estable. Valores propios $-0.25\pm j3.12$.
  \item Arriba, $\cos\pi=-1$: $s^2+0.5s-9.81$, un signo menos: inestable. Valores propios $2.89$ y $-3.39$.
\end{itemize}
Sin fricción ($b=0$), abajo quedaría $s^2+9.81$, con valores propios sobre el eje: la linealización no decidiría (aunque, por conservación de la energía, el péndulo sin fricción es marginalmente estable).
\end{ejemplo}

\begin{pausa}
¿Para qué valores de $a$ es Hurwitz $\Am(a)=\begin{bmatrix}0&1&0\\0&0&1\\-a&-3&-2\end{bmatrix}$? ¿A qué frecuencia oscila en el límite?
\respuesta{$\det(s\Id-\Am)=s^3+2s^2+3s+a$. Por \eqref{eq:es-tercer}: $0<a<6$. En $a=6$ el auxiliar es $2s^2+6$: $\omega=\sqrt3=1.732$\,rad/s.}
\end{pausa}

%% ---------------------------------------------------------------------
\section{Robustez: el teorema de Kharitonov}
\label{sec:es-kharitonov}
%% ---------------------------------------------------------------------

Routh-Hurwitz decide la estabilidad de \emph{un} polinomio. Pero los parámetros de una planta real nunca se conocen exactamente: una ganancia está entre $0.8$ y $1.3$, un amortiguamiento entre $0.2$ y $0.3$, una frecuencia natural entre $10$ y $12$\,rad/s. El modelo es en realidad una \emph{familia} de modelos, y lo que queremos es que el lazo sea estable para \emph{todos} ellos: que sea \clave{robustamente estable}.

Con la planta de segundo orden de ese ejemplo, $\Gp(s)=\dfrac{\Kproc}{s^2+2\zeta\wn s+\wn^2}=\dfrac{q_0}{r_0s^2+r_1s+r_2}$, los intervalos son $q_0\in[0.8,\,1.3]$, $r_0=1$, $r_1=2\zeta\wn\in[4,\,7.2]$ y $r_2=\wn^2\in[100,\,144]$. Una planta cuyos coeficientes varían en intervalos se llama \clave{planta de intervalo}. Con un controlador $\Gc=N_c/D_c$, la ecuación característica $D_cD_p+N_cN_p=0$ tiene coeficientes que también varían en intervalos, y hay infinitos polinomios que revisar.

\begin{definicion}[label=def:es-intervalo]{Familia de polinomios de intervalo}
Es el conjunto de polinomios
\[
  \chi(s)=a_0s^n+a_1s^{n-1}+\cdots+a_{n-1}s+a_n,\qquad a_k\in[a_k^-,\,a_k^+],
\]
en el que cada coeficiente varía \emph{independientemente} de los demás en su intervalo (con $0\notin[a_0^-,a_0^+]$). La familia es \clave{robustamente estable} si todos sus polinomios son Hurwitz.
\end{definicion}

\begin{definicion}[label=def:es-kharitonov]{Teorema de Kharitonov}
Una familia de polinomios de intervalo es robustamente estable \alerta{si y solo si} son Hurwitz los cuatro polinomios
\[
\begin{aligned}
  K_1(s)&=a_n^-+a_{n-1}^-\,s+a_{n-2}^+\,s^2+a_{n-3}^+\,s^3+a_{n-4}^-\,s^4+a_{n-5}^-\,s^5+\cdots\\
  K_2(s)&=a_n^++a_{n-1}^+\,s+a_{n-2}^-\,s^2+a_{n-3}^-\,s^3+a_{n-4}^+\,s^4+a_{n-5}^+\,s^5+\cdots\\
  K_3(s)&=a_n^-+a_{n-1}^+\,s+a_{n-2}^+\,s^2+a_{n-3}^-\,s^3+a_{n-4}^-\,s^4+a_{n-5}^+\,s^5+\cdots\\
  K_4(s)&=a_n^++a_{n-1}^-\,s+a_{n-2}^-\,s^2+a_{n-3}^+\,s^3+a_{n-4}^+\,s^4+a_{n-5}^-\,s^5+\cdots
\end{aligned}
\]
escritos desde el término independiente hacia arriba, con los patrones de extremos $(-,-,+,+)$, $(+,+,-,-)$, $(-,+,+,-)$ y $(+,-,-,+)$ que se repiten cada cuatro términos~\citep{Kha78}.
\end{definicion}

Infinitas verificaciones se reducen a cuatro, y a cada una se le aplica Routh-Hurwitz. ¿Por qué esos cuatro y no, simplemente, ``todos los coeficientes en el mínimo'' y ``todos en el máximo''? Porque en $s=j\omega$ el polinomio vale
\[
  \chi(j\omega)=\underbrace{(a_n-a_{n-2}\omega^2+a_{n-4}\omega^4-\cdots)}_{\text{parte real: coeficientes pares}}
  +j\,\omega\underbrace{(a_{n-1}-a_{n-3}\omega^2+a_{n-5}\omega^4-\cdots)}_{\text{parte imaginaria: coeficientes impares}} ,
\]
con signos que se alternan de dos en dos. La parte real es mínima con $(a_n^-,a_{n-2}^+,a_{n-4}^-,\dots)$ y máxima con los extremos opuestos, y lo mismo pasa con la parte imaginaria. Para cada $\omega$, los valores de toda la familia llenan un \emph{rectángulo} del plano complejo cuyos vértices son exactamente $K_1,\dots,K_4$ (\cref{fig:es-rectangulos}). La familia pierde la estabilidad cuando algún polinomio tiene una raíz en el eje imaginario, es decir, cuando algún rectángulo toca el origen; y eso no puede pasar sin que uno de los vértices lo cruce antes.

\begin{figure}[ht]
\centering
\includegraphics[width=0.85\linewidth]{cap05/f_kharitonov}
\caption{Valores de $\chi(j\omega)$ para la familia $[1,2]s^3+[4,5]s^2+[3,4]s+[1,2]$: para cada $\omega$ forman un rectángulo cuyos vértices (puntos de colores) son los cuatro polinomios de Kharitonov. Ningún rectángulo toca el origen: la familia es robustamente estable.}
\label{fig:es-rectangulos}
\end{figure}

\begin{ejemplo}[label=ej:es-khar1]{Un compensador sobre una planta de intervalo}
Sean $\Gc(s)=\dfrac{s+2}{s^2+2s+2}$ y $\Gp(s)=\dfrac{q_0}{r_0s^3+r_1s^2+r_2s+r_3}$, con $r_0\in[1,\,1.1]$, $r_1\in[4,\,4.2]$, $r_2\in[6,\,8]$, $r_3\in[10,\,20]$ y $q_0\in[3,\,5]$. ¿Es robustamente estable el lazo?
\solucion La ecuación característica $(s^2+2s+2)(r_0s^3+r_1s^2+r_2s+r_3)+q_0(s+2)=0$ es
\[
  a_0s^5+a_1s^4+a_2s^3+a_3s^2+a_4s+a_5=0,
\]
con
\[
\begin{aligned}
  a_0&=r_0\in[1,\,1.1], & a_1&=2r_0+r_1\in[6,\,6.4], & a_2&=2r_0+2r_1+r_2\in[16,\,18.6],\\
  a_3&=2r_1+2r_2+r_3\in[30,\,44.4], & a_4&=2r_2+2r_3+q_0\in[35,\,61], & a_5&=2r_3+2q_0\in[26,\,50].
\end{aligned}
\]
Los cuatro polinomios de Kharitonov son (de $s^0$ a $s^5$)
\[
\begin{aligned}
  K_1&: (26,\,35,\,44.4,\,18.6,\,6,\,1) &&\text{Hurwitz}\\
  K_2&: (50,\,61,\,30,\,16,\,6.4,\,1.1) &&\text{dos raíces en }0.364\pm j2.01\\
  K_3&: (26,\,61,\,44.4,\,16,\,6,\,1.1) &&\text{dos raíces en }0.041\pm j2.60\\
  K_4&: (50,\,35,\,30,\,18.6,\,6.4,\,1) &&\text{dos raíces en }0.194\pm j1.50
\end{aligned}
\]
Tres de los cuatro son inestables: la familia de intervalo \alerta{no es robustamente estable}. Pero hay que tener cuidado con la conclusión: los seis coeficientes dependen de solo cinco parámetros, así que no varían independientemente, y la ``caja'' de coeficientes contiene polinomios que la planta real no puede producir. ¿Hay algún caso \emph{real} inestable? Sí: con $r_0=1$, $r_1=4$, $r_2=6$, $r_3=20$ y $q_0=5$, el lazo tiene polos en $0.033\pm j2.26$. El lazo con la planta nominal (valores centrales) sí es estable, pero el diseño no es robusto.
\end{ejemplo}

\begin{cuidado}[Cuando los coeficientes no son independientes]
El teorema supone que cada coeficiente varía por su cuenta. Si varios coeficientes dependen de los mismos parámetros físicos, Kharitonov se aplica a una familia \emph{más grande} que la real: si dice ``robustamente estable'', la conclusión vale (la familia real está adentro); si dice ``no'', el resultado es \alerta{conservador} y hay que buscar un caso real inestable para confirmarlo.
\end{cuidado}

\begin{ejemplo}[label=ej:es-khar-diseno]{Rango robusto de ganancia}
En el lazo de la \cref{fig:es-khar}, $r_1\in[6,\,7]$, $r_2\in[100,\,120]$ y $r_3\in[1,\,2]$. Halle el rango de $K$ para el cual el lazo es robustamente estable.
\solucion La ecuación característica es $(s^2+r_1s+r_2)(r_3s+1)+K=0$:
\[
  r_3s^3+(r_1r_3+1)\,s^2+(r_2r_3+r_1)\,s+(r_2+K)=0,
\]
con $a_0\in[1,\,2]$, $a_1\in[7,\,15]$, $a_2\in[106,\,247]$ y $a_3\in[100+K,\,120+K]$. Para tercer orden, la condición \eqref{eq:es-tercer} es $a_1a_2>a_0a_3$, y el caso más desfavorable es el que hace el lado izquierdo mínimo y el derecho máximo: $a_1^-a_2^-$ contra $a_0^+a_3^+$. Ese es precisamente $K_4$ (los otros tres polinomios son más holgados), y da
\[
  7\cdot106=742>2\,(120+K)\;\Longrightarrow\;K<251 .
\]
Con $a_3>0$ para todos los casos, $K>-100$: Kharitonov garantiza $\alerta{-100<K<251}$.

¿Qué tan conservador es? Aquí $r_3$ aparece en tres coeficientes y $r_1$, $r_2$ en dos. Haciendo la cuenta con los parámetros, la condición es $r_1r_2r_3+r_1^2+r_1/r_3>K$, cuyo mínimo (en $r_1=6$, $r_2=100$, $r_3=1$) es $642$. La familia real tolera $K<642$: Kharitonov da menos de la mitad del rango (\cref{fig:es-khar-diseno}). Con la planta nominal ($r_1=6.5$, $r_2=110$, $r_3=1.5$), Routh daría $K<1119$: diseñar con el modelo nominal sería peligroso.
\end{ejemplo}

\begin{figure}[ht]
\centering
\includegraphics[width=0.72\linewidth]{cap05/lazo_khar}
\caption{Lazo con planta y sensor inciertos del ejemplo~\ref{ej:es-khar-diseno}.}
\label{fig:es-khar}
\end{figure}

\begin{figure}[ht]
\centering
\includegraphics[width=0.85\linewidth]{cap05/f_khar_diseno}
\caption{Ejemplo~\ref{ej:es-khar-diseno}: parte real del polo más a la derecha en el peor caso de la familia real. Kharitonov garantiza $K<251$; el límite real es $642$; con el modelo nominal se creería que es $1119$.}
\label{fig:es-khar-diseno}
\end{figure}

\codigorepo{cap05_estabilidad/C5_Ej7_Kharitonov.m}{kharitonov}{los cuatro polinomios con la función \cmd{kharitonov} y la comparación con un muestreo aleatorio}

\begin{profundiza}[Menos de cuatro polinomios]
Para órdenes bajos no hacen falta los cuatro: con $n=3$ basta revisar uno ($K_4$, el que usamos arriba), con $n=4$ dos y con $n=5$ tres; desde $n=6$ hacen falta los cuatro~\citep{AJM87}. El teorema también tiene límites: exige coeficientes independientes, y no dice nada sobre el desempeño, solo sobre la estabilidad.
\end{profundiza}

\begin{pausa}
La familia $a_3s^3+a_2s^2+a_1s+a_0$ tiene $a_3\in[1,2]$, $a_2\in[4,5]$, $a_1\in[3,4]$ y $a_0\in[1,2]$. ¿Cuántos polinomios hay que verificar? ¿Es robustamente estable?
\respuesta{Cuatro (y en tercer orden basta el peor, $K_4$). Peor caso de \eqref{eq:es-tercer}: $a_2^-a_1^-=4\cdot3=12>a_3^+a_0^+=2\cdot2=4$. Sí es robustamente estable.}
\end{pausa}

%% ---------------------------------------------------------------------
\section{Para llevar}
\label{sec:es-cierre}
%% ---------------------------------------------------------------------

\begin{parallevar}
\begin{enumerate}
  \item La estabilidad es una propiedad del \emph{equilibrio}. Asintóticamente estable: la respuesta natural tiende a cero (todos los polos con parte real negativa). Marginal: se mantiene acotada (polos simples sobre el eje imaginario). Solo la asintótica garantiza BIBO. Los ceros no intervienen.
  \item La estabilidad del lazo se decide en la ecuación característica $1+\Gma(s)=0$.
  \item Condición necesaria: todos los coeficientes presentes y del mismo signo. Suficiente solo para $n\le2$; para $n=3$, además, $a_1a_2>a_0a_3$.
  \item Routh-Hurwitz: el número de raíces con parte real positiva es el número de cambios de signo en la primera columna. Cero aislado en la primera columna: $\epsilon$. Fila de ceros: polinomio auxiliar (raíces simétricas, que se pueden calcular).
  \item Con parámetros en los coeficientes, la tabla da rangos y regiones de diseño. La fila que se anula da $\Kcr$ y el auxiliar, $\wcr$: los mismos números de Ziegler-Nichols y del margen de ganancia. Con $s=z-\sigma$ se exige un margen de estabilidad.
  \item En el espacio de estado se analiza $\det(s\Id-\Am)$, no una función de transferencia simplificada: puede esconder modos inestables. Linealizando, $\Am$ Hurwitz implica estabilidad local del equilibrio no lineal.
  \item Kharitonov: una familia de intervalo es robustamente estable si y solo si cuatro polinomios lo son. Si los coeficientes dependen de los mismos parámetros, el resultado es conservador.
\end{enumerate}
\end{parallevar}

\begin{profundiza}[Correcciones respecto a la versión anterior del capítulo]
(i) Se retiró la sección de tiempo muerto (aproximación de Padé) y sus dos ejercicios: el curso no cubre retardos. (ii) ``Estable si todos los polos y ceros están en el semiplano izquierdo'' $\to$ solo los polos. (iii) $s^3+2s^2+s+2$ es un caso de fila de ceros, no de $\epsilon$; el auxiliar da las raíces $\pm j$. (iv) Se resolvieron los ejemplos que no tenían solución ($s^5+7s^4+6s^3+42s^2+8s+56$, el primer ejemplo de Kharitonov y el diseño robusto). (v) En el primer ejemplo de Kharitonov faltaba el intervalo de $a_4\in[35,61]$, y la familia \emph{no} es robustamente estable. (vi) El rango robusto del diseño es $-100<K<251$ por Kharitonov y $K<642$ con la dependencia entre coeficientes. (vii) $C(s)/R(s)\to Y(s)/R(s)$; Kharitanov $\to$ Kharitonov. (viii) Historia: Hurwitz publicó en 1895 en \emph{Mathematische Annalen}; el premio Adams de Routh fue en 1877; la anécdota del sirviente es de lord Kelvin. Todos los números se verificaron con \texttt{herramientas/cap05\_numeros.py} y \cmd{verificarCap05.m}.
\end{profundiza}

%% ---------------------------------------------------------------------
\section{Problemas}
\label{sec:es-problemas}
%% ---------------------------------------------------------------------

Los problemas marcados con \etiquetaM{} requieren \matlab/\simulink; los marcados con \etiquetaEstrella{} son de profundización.

\subsection*{El criterio}

\begin{problema}[tablas]\label{prob:es-tablas}
Construya la tabla de Routh, cuente las raíces con parte real positiva y clasifique el equilibrio (asintóticamente estable, marginal o inestable): (a) $s^4+3s^3+3s^2+2s+1$; (b) $s^5+2s^4+3s^3+6s^2+2s+1$; (c) $s^5+3s^4+5s^3+15s^2+4s+12$; (d) $s^4+2s^3-3s^2+4s+5$. En (d) justifique \emph{sin} construir la tabla completa. En (c) identifique el polinomio auxiliar y calcule las raíces que están sobre el eje imaginario. \etiquetaM{} Verifique con \cmd{tablaRouth} y \cmd{roots}.
\end{problema}

\begin{problema}[ganancia crítica]\label{prob:es-critica}
Con retroalimentación unitaria negativa: (i) $\Gma(s)=\dfrac{K}{s(s+2)(s+5)}$; (ii) $\Gma(s)=\dfrac{K(s+1)}{s(s-1)(s+6)}$. (a) Escriba las ecuaciones características. (b) Halle el rango de $K$ para estabilidad asintótica. (c) Para (i), halle $\Kcr$ y la frecuencia de la oscilación sostenida. (d) \etiquetaM{} Simule (i) ante un escalón con $K=0.5\Kcr$, $\Kcr$ y $1.5\Kcr$ y mida el periodo en el caso crítico. (e) Explique por qué en (ii) aparece una cota \emph{inferior} para $K$ y verifíquelo ubicando los polos para tres valores de $K$.
\end{problema}

\begin{problema}[ganancias negativas y Bode]\label{prob:es-bode}
Para $\Gma(s)=\dfrac{K}{(s+1)(s+3)(s+5)}$, con $K\in\mathbb{R}$: (a) halle el rango de $K$ (incluidos los negativos); (b) $\Kcr>0$ y la frecuencia de oscilación; (c) ¿cuántos polos en el semiplano derecho hay con $K=-20$? (d) \etiquetaM{} Con $K=1$, compare $\Kcr$ y $\wcr$ con lo que da \cmd{margin}. ¿Por qué el margen de ganancia no le dice nada sobre la cota negativa?
\end{problema}

\subsection*{Diseño}

\begin{problema}[región de un PI]\label{prob:es-pi}
La planta $\Gp(s)=\dfrac{1}{(s+1)(s+2)}$ se controla con $\Gc(s)=\Kc+K_i/s$ y retroalimentación unitaria. (a) Halle las condiciones sobre $\Kc$ y $K_i$ para estabilidad asintótica. (b) Dibuje la región en el plano $(\Kc,K_i)$. (c) \etiquetaM{} Escoja un punto adentro, uno sobre la frontera y uno afuera, y simule. (d) Sobre la frontera, calcule la frecuencia de la oscilación con el auxiliar y compárela con la simulación. (e) Exija que todos los polos tengan parte real menor o igual a $-0.5$ con $s=z-0.5$, y superponga la nueva región. (f) \etiquetaM{} Tome un punto de la región de (b) que no esté en la de (e), simule y explique qué hace inaceptable el diseño.
\end{problema}

\begin{problema}[estabilidad local]\label{prob:es-local}
El sistema $\dot x_1=x_2$, $\dot x_2=-2x_1+x_1^2-c\,x_2$ tiene equilibrios en $x^*=(0,0)$ y $x^*=(2,0)$. (a) Linealice alrededor de cada uno. (b) Con $c=1$, decida la estabilidad de cada equilibrio con Routh-Hurwitz y con los valores propios. (c) ¿Qué dice la linealización si $c=0$? (d) \etiquetaM{} Simule el sistema no lineal con $c=1$ desde $(0.5,0)$ y desde $(2.1,0)$ y relacione lo que ve con (b).
\end{problema}

\subsection*{Espacio de estado}

\begin{problema}[matrices Hurwitz]\label{prob:es-matrices}
(a) Decida si $\Am_1=\begin{bmatrix}0&1&0\\0&0&1\\-6&-11&-6\end{bmatrix}$ y $\Am_2=\begin{bmatrix}-1&2\\3&-4\end{bmatrix}$ son Hurwitz, con Routh sobre $\det(s\Id-\Am)$ y con los valores propios. (b) Halle el rango de $a$ para el que $\Am_3(a)=\begin{bmatrix}0&1&0\\0&0&1\\-a&-3&-2\end{bmatrix}$ es Hurwitz. (c) \etiquetaM{} En el límite superior, calcule la frecuencia de oscilación con el auxiliar y verifíquela simulando desde $x(0)=[1\;\;0\;\;0]^\top$. (d) Con $\Am_1$, $\Bm=[0\;\;0\;\;1]^\top$ y $\Dm=0$, encuentre un $\Cm$ para el que la función de transferencia tenga \emph{menos} polos que valores propios tiene $\Am_1$. ¿Por qué Routh sobre esa función de transferencia no basta para hablar del equilibrio?
\end{problema}

\subsection*{Robustez}

\begin{problema}[familia de cuarto orden]\label{prob:es-familia}
Sea $\chi(s)=a_0s^4+a_1s^3+a_2s^2+a_3s+a_4$ con $a_0\in[1,1.2]$, $a_1\in[4,5]$, $a_2\in[6,8]$, $a_3\in[3,4]$, $a_4\in[1,2]$. (a) Escriba los cuatro polinomios de Kharitonov. (b) Aplique Routh a cada uno y concluya. (c) \etiquetaM{} Genere 2000 polinomios aleatorios de la familia y verifique que ninguno es inestable; compare el costo con (b). (d) Repita (a)--(b) con $a_3\in[0.2,4]$. (e) Si la respuesta a (d) es negativa, muestre un polinomio inestable de la familia y sus raíces.
\end{problema}

\begin{problema}[planta incierta]\label{prob:es-incierta}
\etiquetaEstrella{} Con $\Gc=K$ y $\Gp(s)=\dfrac{1}{s^3+r_1s^2+r_2s+r_3}$, $r_1\in[6,7]$, $r_2\in[100,120]$, $r_3\in[1,2]$: (a) escriba el polinomio característico y los intervalos de sus coeficientes; (b) con Kharitonov, halle el rango robusto de $K$ y diga qué polinomio impone la restricción; (c) compare con el rango de la planta nominal ($6.5$, $110$, $1.5$); (d) escoja un $K$ admisible para la nominal pero no para la familia y muestre un caso inestable; (e) \etiquetaM{} recorra los ocho vértices $(r_1,r_2,r_3)$ y verifique el peor caso. (f) ¿Es conservador aquí el resultado de Kharitonov? Compare con el ejemplo~\ref{ej:es-khar-diseno}.
\end{problema}

%% ---------------------------------------------------------------------
\section{Soluciones seleccionadas}
%% ---------------------------------------------------------------------

\paragraph{Problema~\ref{prob:es-tablas}}
(a) Sin cambios: asintóticamente estable. (b) Un cero en la fila $s^3$ ($\epsilon$) y dos cambios: inestable, raíces $0.096\pm j1.62$. (c) Fila $s^3$ de ceros; auxiliar $3s^4+15s^2+12=3(s^2+1)(s^2+4)$: raíces $\pm j$ y $\pm2j$, sin cambios: marginalmente estable ($s^5+\dots=(s+3)(s^2+1)(s^2+4)$). (d) Hay un coeficiente negativo: no es estable (la tabla da dos cambios: $0.956\pm j1.115$).

\paragraph{Problema~\ref{prob:es-critica}}
(i) $s^3+7s^2+10s+K$: $0<K<70$; $\Kcr=70$, auxiliar $7s^2+70$, $\wcr=\sqrt{10}=3.16$\,rad/s ($\Pcr=1.99$\,s). (ii) $s^3+5s^2+(K-6)s+K$: $K>0$ y $\frac{5(K-6)-K}{5}>0$, o sea $K>7.5$. Con poca ganancia el polo en $+1$ no alcanza a moverse al semiplano izquierdo.

\paragraph{Problema~\ref{prob:es-bode}}
(a) $s^3+9s^2+23s+15+K$: $-15<K<192$. (b) $\Kcr=192$, $\wcr=\sqrt{23}=4.80$\,rad/s. (c) Con $K=-20$ el término independiente es $-5$: un cambio de signo, un polo en $+0.20$. (d) \cmd{margin} da $\GM=192$ en $4.80$\,rad/s; la cota negativa corresponde a un polo que cruza por el origen ($\omega=0$), que el margen de ganancia (definido en el cruce de $-180^\circ$) no mira.

\paragraph{Problema~\ref{prob:es-pi}}
(a) $s^3+3s^2+(2+\Kc)s+K_i$: $K_i>0$ y $K_i<3(2+\Kc)$. (d) Sobre la frontera, el auxiliar $3s^2+K_i$ da $\omega=\sqrt{2+\Kc}$. (e) $z^3+1.5z^2+(\Kc-0.25)z+(K_i-\Kc/2-0.375)$: $\frac{\Kc}{2}+0.375<K_i<2\Kc$. (f) Por ejemplo $(\Kc,K_i)=(2,5)$: estable, pero con polos en $-0.39\pm j1.45$, poco amortiguados.

\paragraph{Problema~\ref{prob:es-local}}
(a) $\Am=\begin{bmatrix}0&1\\-2+2x_1^*&-c\end{bmatrix}$. (b) En $(0,0)$: $s^2+s+2$, estable ($-0.5\pm j1.32$). En $(2,0)$: $s^2+s-2$, inestable ($1$ y $-2$). (c) Con $c=0$, en el origen los valores propios son $\pm j\sqrt2$: la linealización no decide.

\paragraph{Problema~\ref{prob:es-matrices}}
(a) $\Am_1$: $s^3+6s^2+11s+6$, Hurwitz. $\Am_2$: $s^2+5s-2$, no ($0.372$ y $-5.372$). (b) $0<a<6$. (c) $\omega=\sqrt3$. (d) Por ejemplo $\Cm=[2\;\;3\;\;1]$: el numerador $s^2+3s+2=(s+1)(s+2)$ cancela dos polos y queda $1/(s+3)$.

\paragraph{Problema~\ref{prob:es-familia}}
(b) Los cuatro son Hurwitz: robustamente estable. (d) Con $a_3\in[0.2,4]$, $K_1$ y $K_4$ son inestables. (e) $K_1$: $s^4+5s^3+8s^2+0.2s+1$, con raíces $0.026\pm j0.350$.

\paragraph{Problema~\ref{prob:es-incierta}}
(a) $s^3+r_1s^2+r_2s+(r_3+K)$, con coeficientes independientes. (b) Peor caso de \eqref{eq:es-tercer}: $6\cdot100>2+K$: $-1<K<598$. (c) Nominal: $-1.5<K<713.5$. (f) No es conservador: cada coeficiente depende de un solo parámetro, y el peor vértice es una planta real.

\respuestaspausas

